\section{The exact bridge and its inverse}
\label{app:bridge-exact}

\paragraph{Forward map.}
With conventions as in appendix~\ref{app:conventions},
\begin{equation}
\Phi(F)_{ab} = \frac{1}{5!}\,F_{\mu_1\dots\mu_5}\,
              (C\Gamma^{\mu_1\dots\mu_5})_{ab},
\qquad a,b = 1,\dots,\spinorDim .
\end{equation}
In the implementation $F$ is carried as a vector of $\dimSelfDualModule$
independent components in a fixed ordering of index five-tuples, and $\Phi$ is a
$\dimSymSpinor \times \dimSelfDualModule$ integer matrix reduced mod $p$.

\paragraph{Gamma trace.}
The gamma trace $(\gamma\text{-}\mathrm{tr}\,\psi)_\mu
 = (\Gamma^\nu)^{ab}\psi_{ab}\,\eta_{\mu\nu}$-type contraction has rank
$\gammaTraceRank$ on $\mathrm{Sym}^2 S_+$. Hence
\begin{equation}
\dim \ker(\gamma\text{-}\mathrm{tr})
 = \dimSymSpinor - \gammaTraceRank
 = \dimGammaTraceless
 = \dim \Lambda^5_+ .
\end{equation}
The certified statement is that $\mathrm{im}\,\Phi \subseteq
\ker(\gamma\text{-}\mathrm{tr})$ and that the inclusion is an equality by rank.

\paragraph{Rank and kernel.}
$\mathrm{rank}\,\Phi = \bridgeForwardRank$ on $\Lambda^5_+$, and
$\ker \Phi \cap \Lambda^5_+ = 0$. Both are computed by exact elimination over
$\Fp$ at each prime used. On $\Lambda^5_-$ the map vanishes identically, which is
the statement that appendix~\ref{app:orientation} shows can silently invert.

\paragraph{Left inverse and round trip.}
Because $\Phi$ has full column rank, a left inverse $\Phi^-$ exists over $\Fp$
and is constructed explicitly by exact elimination. The certified identity is
\begin{equation}
\Phi^- \Phi = \mathrm{id}_{\dimSelfDualModule},
\end{equation}
checked as a matrix identity, not merely on sampled vectors. Note that
$\Phi\Phi^-$ is \emph{not} the identity on $\mathrm{Sym}^2 S_+$; it is the
projector onto the gamma-traceless subspace, which is the correct expectation and
is checked as such.

\paragraph{Equivariance.}
For $\Lambda$ in the Lorentz group acting on the tensor side and $S(\Lambda)$ the
corresponding spinor representative, the certified statement is
\begin{equation}
\Phi(\Lambda \cdot F) = \chi(\Lambda)\,
 S(\Lambda)\,\Phi(F)\,S(\Lambda)^{T}
\end{equation}
where the character $\chi$ is solved for rather than assumed. In practice
$\chi \equiv 1$ in the normalisation of appendix~\ref{app:conventions}; solving
for it rather than assuming it is what turns the equivariance check into a test
that can fail.

\paragraph{Normalisation independence.}
Rescaling $\Phi \to \lambda\Phi$ rescales $\Phi^- \to \lambda^{-1}\Phi^-$ and
leaves every rank, dimension and intersection in this paper unchanged. The
normalisation is therefore fixed for tidiness only.

\paragraph{Test inventory.}
The bridge suite contains $\nBridgeTests$ tests covering: the Clifford relations;
$\star^2 = +\starSquared$; the symmetry of $C\Gamma^{(5)}$; the gamma-trace rank;
the forward rank; the vanishing on the opposite eigenspace; the left inverse; the
round trip; the projector property of $\Phi\Phi^-$; equivariance with the
character solved for; and the orientation tests of
appendix~\ref{app:orientation}, including the deliberate branch flip that must
break.
